\documentclass[12pt]{article}
\usepackage{amsmath,amssymb}

\begin{document}
\section*{{2025 AMC 12B Problems/Problem 23}}
Source: AoPS Wiki

\subsection*{Solution}
Notice that there are $z$ distinct remainders modulo $z$ . However, if we let $R(x)$ denote the remainder when $2025x$ is divided by $z$ , we know that by the problem condition, \[R(0)<R(1)<R(2)<\cdots<R(z-1)\] But there are only $z$ distinct remainders, and each of the $z$ terms above must correspond to a distinct remainder, so we must have $R(k)=k$ for all $k=0,1,\dots,z-1$ . Then $2025k\equiv k\pmod{z}$ , so $2024k\equiv 0\pmod{z}$ . Since $k$ can vary, this is equivalent to $2024\equiv 0\pmod{z}$ , or $z\mid 2024$ . Therefore, the values $z$ that satisfy the property are the factors of $2024$ , so we simply need to find the sum of the factors of $2024$ and subtract $1$ at the end since $z\ne 1$ . But $2024=2^3\cdot 11\cdot 23$ , so the sum of the factors minus $1$ is \[(1+2+4+8)(1+11)(1+23)-1=15\cdot12\cdot24-1=4320-1=\boxed{\textbf{(E)}~4319 }.\] ~ eevee9406

\end{document}
